\documentclass[11pt,a4paper]{article}
\usepackage[margin=1in]{geometry}
\usepackage[utf8]{inputenc}
\usepackage[english]{babel}
\usepackage{amsmath,amssymb}
\usepackage{microtype}
\usepackage{hyperref}

\setlength{\parindent}{0pt}
\setlength{\parskip}{0.75\baselineskip}
\sloppy

\title{A Formalization Program for Field 01\\\large Phase Nodes, Scalar Profiles, Normal-Retention Interpretation, Memory, and Boundary Recording}
\author{Nikolai Dereviankin}
\date{}

\newcommand{\field}{\mathcal{F}_{01}}
\newcommand{\phase}{\varphi}
\newcommand{\normal}{\mathcal{N}}
\newcommand{\memory}{\mathcal{M}}
\newcommand{\boundary}{\partial}
\newcommand{\bulk}{\mathrm{bulk}}
\newcommand{\wind}{\mathrm{wind}}
\newcommand{\cov}{\mathrm{cov}}

\begin{document}
\maketitle

\begin{center}
\textit{Research version v0.3 --- English formalization text with standard-first diagnostic appendix}
\end{center}

\begin{abstract}
This paper outlines a cautious formalization program for Field 01. It organizes the model's current mathematical development into a chain of concepts: phase, closed circulation, scalar modulus / VEV-profile layer, labelled normal-retention interpretation, gauge-like compensation, current-like flow, memory as an equivalence class, boundary recording, and reduced external description. The proposal is not presented as a completed physical theory, nor as a replacement for the Standard Model, quantum field theory, general relativity, or black-hole thermodynamics. Instead, it gives a structured toy-level framework in which the book-level language of Field 01 can be compared with known mathematical structures, especially phase/gauge/vortex models. The central result is not a new fundamental equation, but a controlled formal vocabulary: closed phase nodes can be represented by winding, the scalar modulus / VEV-profile layer by a scalar profile, the normal-retention reading by an interpretation placed on that layer, screening by a connection-like field, and memory by equivalence classes of preserved phase-scalar-gauge data. The v0.3 diagnostic layer added here translates selected Field 01 phrases into standard Abelian-Higgs benchmarks, scalar boundary-layer tension, gauge-compensation stiffness, overlap-motivated distance bins, and Wilson-line dressed phase bookkeeping. The horizon interpretation is expressed as a bulk-to-boundary memory map and as reduced external access to boundary data. All claims remain provisional and require further comparison with established physics.
\end{abstract}

\textbf{Keywords:} Field 01; formalization program; phase circulation; closed node; scalar modulus; VEV-profile; normal-retention interpretation; memory; gauge-like compensation; Abelian-Higgs diagnostics; boundary-layer tension; Wilson-line phase; vortex models; boundary recording; reduced density matrix.

\section{How to Read This Paper}

This paper is a formalization roadmap, not a completed theory. It should be read as the third research text in the current Field 01 sequence. The first paper introduced the elementary particle as a closed phase configuration. The second paper interpreted a black-hole horizon as a phase recording surface. The present paper asks a narrower question: what mathematical skeleton is needed for these claims to become discussable and testable?

Three levels must be kept separate:
\begin{enumerate}
\item \textbf{Established mathematics and physics}: phase fields, winding numbers, covariant derivatives, currents, reduced density matrices, vortex-like structures, and black-hole thermodynamic formulas.
\item \textbf{Field 01 interpretation}: phase as rhythm, a scalar modulus / VEV-profile layer as the carrier of a labelled normal-retention reading, memory as preserved phase structure, and horizon recording as a boundary representation of local phase data.
\item \textbf{Open hypotheses}: the identification of the scalar-profile layer, under the normal-retention interpretation, with particle-like mass; the use of boundary memory for horizons; and any relation to charge, spin, or black-hole information.
\end{enumerate}

The aim is not to claim novelty where standard mathematics already exists. In fact, one important conclusion is that the first toy formalization of Field 01 naturally approaches known Abelian-Higgs or vortex-like structures. The possible value of Field 01 lies in its interpretation and in the proposed link between scalar-profile retention, memory, and boundary recording.

\section{Standard-First Vocabulary Rule}

This text follows a standard-first rule. Each technical object should be introduced first with the nearest standard vocabulary, then with the Field 01 reading, and finally with an explicit statement of what remains hypothetical. The reference map for this rule is maintained separately in \texttt{FIELD01\_REFERENCE\_MAP.md}.

\begin{itemize}
\item \textbf{phase rhythm}: first write phase variable, angular variable, or phase of a complex scalar; the rhythm language is Field 01-specific.
\item \textbf{closed phase node}: first write vortex-like winding, topological defect, or winding number; the particle interpretation remains open.
\item \textbf{scalar modulus / VEV-profile layer}: first write scalar amplitude, radial profile, or order-parameter modulus; the normal-retention interpretation is only the Field 01 label placed on this layer.
\item \textbf{gauge-like compensation}: first write connection-like field, covariant derivative, or screening; this is not a derivation of electromagnetism.
\item \textbf{memory}: first write preserved invariant data or equivalence class; the memory interpretation remains hypothetical.
\item \textbf{boundary recording}: first write boundary data, reduced state, or holographic-style boundary language; no solution of black-hole information is claimed.
\end{itemize}

\section{Formalization Chain}

The current formalization can be summarized as
\begin{align}
&\text{phase}
\rightarrow \text{closed node}
\rightarrow \text{scalar modulus / VEV-profile layer}
\nonumber\\
&\rightarrow \text{normal-retention interpretation}
\rightarrow \text{gauge-like compensation}
\rightarrow \text{current}
\nonumber\\
&\rightarrow \text{memory class}
\rightarrow \text{boundary map}
\rightarrow \text{reduced state}.
\end{align}

This chain is not a derivation of known physics. It is a formalization route from the conceptual language of the model to mathematical objects that can be analyzed.

The main variables are
\begin{equation}
\field, \qquad \phase, \qquad N \;\text{or}\; \normal, \qquad A_\mu, \qquad \memory.
\end{equation}

Here $\phase$ denotes a phase variable, $N$ is the scalar toy profile or order-parameter modulus, $\normal$ denotes the optional Field 01 normal-retention reading of that scalar layer, $A_\mu$ is a compensating connection-like field, and $\memory$ denotes preserved phase-structural data. At this stage $\field$ is not yet defined by a final Lagrangian.

\section{Phase And Closed Circulation}

Let
\begin{equation}
\phase:X\to S^1
\end{equation}
be a phase variable on an effective domain $X$. For a contour $\mathcal{C}$, define the phase circulation
\begin{equation}
\Gamma[\mathcal{C}]=\oint_{\mathcal{C}}d\phase.
\end{equation}

A closed phase node is represented by the condition
\begin{equation}
\oint_{\mathcal{C}}d\phase=2\pi n, \qquad n\in\mathbb{Z}.
\end{equation}

Equivalently,
\begin{equation}
Q_{\wind}=\frac{1}{2\pi}\oint_{\mathcal{C}}d\phase=n.
\end{equation}

In the toy model one uses the ansatz
\begin{equation}
\phase_n(r,\theta)=n\theta.
\end{equation}

This gives the desired winding, but a bare winding has long-range energy cost. On an annulus $a\le r\le R$,
\begin{equation}
\int_0^{2\pi}\int_a^R \frac{n^2}{r^2}r\,dr\,d\theta
=2\pi n^2\log\frac{R}{a}.
\end{equation}

Thus a bare phase winding is not yet a localized finite-energy particle-like object.

\section{Scalar Profile, Normal-Retention Interpretation, And Mass Proxy}

In vortex-like and Abelian-Higgs-type models, a winding configuration is usually accompanied by a scalar amplitude or radial profile. Field 01 first uses the same kind of scalar profile as a toy proxy. Only after this standard scalar layer is introduced does the book-level language label it as normal-retention interpretation:
\begin{equation}
N=N(r).
\end{equation}

A minimal unscreened radial energy is
\begin{equation}
E_n[N]
=2\pi\int_0^R
\left[
\frac{A}{2}(N')^2
+\frac{B}{2}\frac{n^2N^2}{r^2}
+\frac{\lambda}{4}(N^2-N_0^2)^2
\right]r\,dr.
\end{equation}

The corresponding radial equation is
\begin{equation}
A\left(N''+\frac{1}{r}N'\right)
-B\frac{n^2}{r^2}N
-\lambda N(N^2-N_0^2)=0.
\end{equation}

In Field 01 language, the mass proxy is the energy of retaining a closed phase-scalar configuration, with the normal-retention interpretation added as a model-specific reading:
\begin{equation}
mc^2 \sim E_{\mathrm{closed}}[\phase,N].
\end{equation}

This is not yet a derivation of physical mass. It is only a formal expression of the idea that mass is associated with the energy cost of maintaining a closed local structure.

A numerical toy check with $A=B=\lambda=N_0=1$, $n=1$, and $r\in[10^{-3},12]$ produced a particle-like profile with finite-disk energy approximately
\begin{equation}
E_{\mathrm{unscreened}}\approx 9.0093.
\end{equation}

A horizon-like boundary condition $N(R_H)=0$ also admits a profile while preserving winding $Q_{\wind}=1$.

\section{Connection-Like Compensation}

In standard gauge notation, a connection field can screen or compensate a long-range phase gradient through a covariant derivative. The Field 01 toy model uses the same structure to represent compensation of retained phase mismatch. Introduce a connection-like field $A_\mu$ and the covariant phase derivative
\begin{equation}
D_\mu\phase=\partial_\mu\phase-A_\mu.
\end{equation}

Under the gauge-like transformation
\begin{equation}
\phase\mapsto\phase+\chi,
\qquad
A_\mu\mapsto A_\mu+\partial_\mu\chi,
\end{equation}
we have
\begin{equation}
D_\mu\phase\mapsto D_\mu\phase.
\end{equation}

The field strength
\begin{equation}
F_{\mu\nu}=\partial_\mu A_\nu-\partial_\nu A_\mu
\end{equation}
is also invariant.

For the radial ansatz, the replacement is
\begin{equation}
\frac{n}{r}\longrightarrow\frac{n-a(r)}{r}.
\end{equation}

A screened toy energy is
\begin{equation}
E_n[N,a]
=2\pi\int_0^R
\left[
\frac{A}{2}(N')^2
+\frac{B}{2}\frac{(n-a)^2N^2}{r^2}
+\frac{C}{2}\frac{(a')^2}{r^2}
+\frac{\lambda}{4}(N^2-N_0^2)^2
\right]r\,dr.
\end{equation}

The radial equations are
\begin{equation}
A\left(N''+\frac{1}{r}N'\right)
-B\frac{(n-a)^2}{r^2}N
-\lambda N(N^2-N_0^2)=0,
\end{equation}

\begin{equation}
a''-\frac{1}{r}a'
+\frac{B}{C}(n-a)N^2=0.
\end{equation}

For the same finite disk, the screened particle-like numerical profile gave
\begin{equation}
E_{\mathrm{screened}}\approx 3.6340,
\end{equation}
which is lower than the unscreened toy energy. This supports the mathematical role of compensation, but it does not yet derive electromagnetism.

\section{Current-Like Object}

A covariant-looking toy density may be written as
\begin{equation}
\mathcal{L}_{\mathrm{toy}}
=\frac{A_N}{2}(\partial_\mu N)(\partial^\mu N)
+\frac{B}{2}N^2(D_\mu\phase)(D^\mu\phase)
-V(N)
-\frac{C}{4}F_{\mu\nu}F^{\mu\nu}.
\end{equation}

The global phase-shift symmetry
\begin{equation}
\phase\mapsto\phase+\epsilon, \qquad \epsilon=\mathrm{const}
\end{equation}
gives the current-like object
\begin{equation}
J^\mu=\frac{\partial\mathcal{L}_{\mathrm{toy}}}{\partial(\partial_\mu\phase)}
=BN^2D^\mu\phase.
\end{equation}

The phase equation gives
\begin{equation}
\partial_\mu J^\mu=0.
\end{equation}

This is a conserved phase-compensation current inside the toy model. It is not yet electric current. It only shows that retained covariant phase flow has a conserved current-like expression.

\section{Relation To Abelian-Higgs And Vortex Models}

The screened toy model is close to known Abelian-Higgs or vortex-like structures. This must be stated explicitly. If one writes
\begin{equation}
\Psi=Ne^{i\phase},
\end{equation}
then
\begin{equation}
|D_\mu\Psi|^2
=(\partial_\mu N)(\partial^\mu N)
+N^2(\partial_\mu\phase-A_\mu)(\partial^\mu\phase-A^\mu).
\end{equation}

Thus the pair $(N,\phase)$ is mathematically close to polar variables of a complex scalar field. Field 01 should not claim this structure as a new invention.

The possible contribution of Field 01 is interpretive:
\begin{itemize}
\item $\phase$ is read as phase rhythm or mode-$1$ circulation;
\item $N$ is first the scalar modulus / VEV-profile layer, then read in Field 01 through the normal-retention interpretation or local-depth label;
\item winding is read as phase memory;
\item $A_\mu$ is read as compensating field response;
\item $N\to0$ is read as the disappearance of local normal depth at a boundary.
\end{itemize}

Whether this interpretation leads to new mathematics or physics remains an open question.

\section{Memory As Equivalence Class}

Equivalence classes and invariant data are standard mathematical tools for separating physical or structural information from choices of representation. The book-level concept of memory is therefore formalized here as an equivalence class of preserved phase-scalar-gauge data, with normal-retention kept as an interpretation of the scalar component:
\begin{equation}
\memory_{\bulk}
=[\phase,N,A_\mu]_{\sim_{\bulk}}.
\end{equation}

A first invariant set is
\begin{equation}
\mathcal{I}_{\bulk}
=(Q_{\wind},\Phi_F,\mathcal{J},E_{\mathrm{class}},\mathcal{B}),
\end{equation}
where $Q_{\wind}$ is winding, $\Phi_F$ denotes flux or connection data, $\mathcal{J}$ denotes current flux, $E_{\mathrm{class}}$ is an energy class, and $\mathcal{B}$ denotes boundary-accessible phase data.

The equivalence relation identifies configurations that preserve these invariants while quotienting out gauge choice, coordinate representation, and small smooth deformations that do not change the record.

Boundary memory is written as
\begin{equation}
\memory_{\boundary}
=[\phase_{\boundary},A_{\boundary},F_{\boundary},Q_{\boundary}]_{\sim_{\boundary}}.
\end{equation}

The bulk-to-boundary memory map is
\begin{equation}
\Pi_{\boundary}:\memory_{\bulk}\to\memory_{\boundary}.
\end{equation}

This is the formal version of the Field 01 statement that local memory becomes boundary phase record.

\section{Boundary Reduction And Thermality}

A standard reduced density matrix is obtained by tracing inaccessible degrees of freedom:
\begin{equation}
\rho_{\mathrm{external}}
=\mathrm{Tr}_{\mathrm{hidden}}\rho_{\mathrm{full}}.
\end{equation}

In Field 01 language, boundary memory may be split as
\begin{equation}
\memory_{\boundary}
=(\memory_{\mathrm{acc}},\memory_{\mathrm{hid}},\mathcal{C}_{\mathrm{corr}}),
\end{equation}
where $\memory_{\mathrm{acc}}$ is accessible boundary data, $\memory_{\mathrm{hid}}$ is inaccessible data, and $\mathcal{C}_{\mathrm{corr}}$ denotes correlations.

Then one may write schematically
\begin{equation}
\rho_{\mathrm{external}}
=\mathrm{Tr}_{\memory_{\mathrm{hid}}}\rho(\memory_{\boundary}).
\end{equation}

A reduced state may appear thermal,
\begin{equation}
\rho_{\mathrm{external}}\approx\frac{e^{-\beta H_{\mathrm{eff}}}}{Z}.
\end{equation}

The Field 01 interpretation is that thermality may reflect reduced access to the full boundary record, not necessarily destruction of the record. This does not derive Hawking radiation or solve the black-hole information problem. It is only a formal expression of the model's interpretive stance.

\section{Technical Appendix: Standard Diagnostics For Current Terminology}

This appendix records the current standard-first diagnostic layer. It does not claim a new particle model, a new force, a new topological invariant, or a solution of the black-hole information problem. Its narrower purpose is to state which Field 01 terms have been translated into standard Abelian-Higgs-type calculations or into explicit relation-class bookkeeping. All Field 01 language in this section remains interpretive unless it is tied to a standard calculated quantity.

\subsection{Standard Vortex Substrate}

Use a complex scalar field and Abelian gauge profile in polar variables,
\begin{equation}
\Psi=N(r)e^{in\theta},
\qquad
A=a(r)\,d\theta.
\end{equation}

The radial energy convention used for the diagnostic benchmark is
\begin{equation}
E
=2\pi\int_0^R dr\left[
\frac{r}{2}(N')^2
+\frac{(n-a)^2N^2}{2r}
+\frac{(a')^2}{2g^2r}
+\frac{\lambda r}{4}(N^2-N_0^2)^2
\right].
\end{equation}

The finite-energy boundary behavior is represented by
\begin{equation}
N(0)=0,
\qquad
a(0)=0,
\qquad
N(R)\approx N_0,
\qquad
a(R)\approx n.
\end{equation}

At the critical point in this normalization,
\begin{equation}
\lambda_{\mathrm{BPS}}=\frac{g^2}{2}.
\end{equation}

For $g=1$, $N_0=1$, $n=1$, and $\lambda=0.5$, the numerical benchmark gives
\begin{equation}
E_{\mathrm{total}}=3.141567,
\qquad
\frac{E_{\mathrm{total}}}{\pi}=0.999992.
\end{equation}

This is standard Abelian-Higgs vortex machinery. Field 01 does not claim this calculation as new. Its role here is to fix the standard substrate before adding any interpretation.

\subsection{Scalar Boundary Diagnostic}

To test the phrase ``loss of retained profile,'' compare the standard outer behavior $N(R)=N_0$ with a forced outer-zero condition $N(R)=0$, while keeping the gauge endpoint $a(R)=n$. Define
\begin{equation}
\Delta E_{\mathrm{outer\ zero}}
=E[N(R)=0,a(R)=n]-E[N(R)=N_0,a(R)=n].
\end{equation}

The raw penalty grows approximately with circumference. The normalized penalty
\begin{equation}
\sigma_R=\frac{\Delta E_{\mathrm{outer\ zero}}}{2\pi R}
\end{equation}
approaches the standard scalar half-kink boundary tension. For $\lambda=1$ and $N_0=1$,
\begin{equation}
\sigma_R\to \frac{2}{3\sqrt{2}}\approx0.471404521.
\end{equation}

More generally, the flat-boundary scaling is
\begin{equation}
\sigma_\infty(\lambda,N_0)
=\frac{2}{3\sqrt{2}}\sqrt{\lambda}\,N_0^3.
\end{equation}

Thus, in this diagnostic, the Field 01 phrase ``loss of retained profile'' reduces to standard scalar boundary-layer physics. This does not prove a new invariant, a particle mass, or a new boundary law.

\subsection{Gauge-Compensation Diagnostics}

The first gauge-compensation diagnostic keeps $N(R)=N_0$ and varies the endpoint gauge profile $a(R)=\alpha$. Define
\begin{equation}
\Delta E_a(R,\alpha)
=E[N(R)=N_0,a(R)=\alpha]-E[N(R)=N_0,a(R)=n].
\end{equation}

For the finite-disk diagnostic, the leading behavior is
\begin{equation}
\Delta E_a(R,\alpha)
\approx \frac{\pi}{gR}(n-\alpha)^2.
\end{equation}

Endpoint failure of gauge compensation is therefore read here as a standard finite-boundary mismatch cost.

The second diagnostic keeps the endpoint flux fixed and perturbs only the interior profile,
\begin{equation}
a_\varepsilon(r)=a_0(r)+\varepsilon f(r),
\qquad
f(0)=f(R)=0.
\end{equation}

The leading cost is the standard second variation,
\begin{equation}
\Delta E[\varepsilon f]
=\varepsilon^2 K[f]+O(\varepsilon^3),
\end{equation}
with
\begin{equation}
K[f]
=2\pi\int_0^R dr\left[
\frac{N^2f^2}{2r}
+\frac{(f')^2}{2g^2r}
\right].
\end{equation}

This identifies compensation stiffness with the positive Hessian / second variation around the standard vortex solution. It does not define a new compensation invariant or a new interaction law.

\subsection{Two-Node Relation As Resolution-Class Bookkeeping}

A two-node retained relation is currently represented as a resolution-dependent equivalence class, not as a force. A gauge-aware relation class may be written schematically as
\begin{equation}
I_{\mathrm{rel}}^{R}(A,B)=
\left(
W_{\mathrm{pair}},
D_{\mathrm{overlap}}(A,B),
\beta_{\mathrm{regime}},
P_q^\gamma(A,B),
B_{\mathrm{pair}}(A,B),
\gamma_{\mathrm{class}}(A,B)
\right).
\end{equation}

Here $W_{\mathrm{pair}}$ is a winding-pair label, $D_{\mathrm{overlap}}$ is a distance/range bin motivated by profile overlap, $\beta_{\mathrm{regime}}$ is a standard type-I/BPS/type-II regime label, $P_q^\gamma$ is a Wilson-line dressed phase bin, $B_{\mathrm{pair}}$ is a boundary-pair label, and $\gamma_{\mathrm{class}}$ records the connector or access-path class. In this reading, two configurations have the same retained relation only when the selected class data match under the declared resolution package $R$.

To reduce arbitrariness in the distance bins, define scalar and gauge tails
\begin{equation}
u_N(r)=1-N(r),
\qquad
u_a(r)=1-a(r),
\end{equation}
and the combined normalized overlap proxy
\begin{equation}
O_{Na}(d)=
\frac{\int d^2x\,[u_N^A u_N^B+u_a^A u_a^B]}
{\int d^2x\,[u_N^2+u_a^2]}.
\end{equation}

For the current standard profile, the combined-overlap thresholds are
\begin{equation}
O_{Na}=0.10 \quad\Rightarrow\quad d=5.400085,
\qquad
O_{Na}=0.01 \quad\Rightarrow\quad d=8.642013.
\end{equation}

A simple mass-weighted tail proxy gives close thresholds,
\begin{equation}
J_{\mathrm{mw}}=0.10 \quad\Rightarrow\quad d=5.183276,
\qquad
J_{\mathrm{mw}}=0.01 \quad\Rightarrow\quad d=8.440450.
\end{equation}

This makes the distance bins less arbitrary, but it is still not a true two-vortex interaction energy. The bins remain diagnostic conventions unless a physical observable and cutoff rule are justified.

\subsection{Gauge-Aware Phase Relation}

A raw phase bin based on $\theta_B-\theta_A$ is gauge-sensitive. The safer standard object is the Wilson-line dressed phase along a connector $\gamma$:
\begin{equation}
h_\gamma(A,B)
=\left[(\theta_B-\theta_A)-\int_\gamma A\right]\bmod 2\pi.
\end{equation}

The corresponding phase bin is
\begin{equation}
P_q^\gamma(A,B)
=\left\lfloor \frac{q h_\gamma(A,B)}{2\pi}\right\rfloor.
\end{equation}

Because an open Wilson-line phase depends on the chosen connector unless a path/access rule is specified, the relation class must include $\gamma_{\mathrm{class}}$ or an explicit boundary-selected connector. The Field 01 reading is therefore not a raw relative phase, but a gauge-aware holonomy-dressed phase class under a declared access rule. This is not entanglement, not a force law, and not a new invariant beyond standard gauge-theory bookkeeping.

\subsection{Current Diagnostic Status}

The current technical status is deliberately modest. Several Field 01 phrases now have standard mathematical diagnostics attached to them:
\begin{itemize}
\item closed field object $\rightarrow$ Abelian-Higgs radial vortex ansatz;
\item retained local profile $\rightarrow$ scalar modulus $N(r)$;
\item loss of retention $\rightarrow$ scalar boundary layer / half-kink tension;
\item phase compensation $\rightarrow$ gauge profile $a(r)$ and covariant derivative;
\item compensation stiffness $\rightarrow$ second variation around a vortex profile;
\item retained relation $\rightarrow$ resolution-dependent equivalence class;
\item distance relation $\rightarrow$ overlap-motivated class bin;
\item phase relation $\rightarrow$ Wilson-line dressed holonomy bin.
\end{itemize}

This is progress in vocabulary control, not a claim of new physics. The next burden is to test whether any of these diagnostics leads to a genuinely new, physically distinguishable structure rather than to a restatement of known Abelian-Higgs, vortex, gauge, or boundary-layer objects.

\section{Immediate Reference Anchors}

The following standard literature should be used as comparison anchors. These references do not support the Field 01 interpretation by themselves; they identify nearby established structures that the text must not claim as new.

\begin{itemize}
\item Nielsen and Olesen, ``Vortex-line models for dual strings,'' \textit{Nuclear Physics B} 61, 45--61 (1973), for the classic scalar-field, winding, and gauge-vortex comparison.
\item Abrikosov, ``On the magnetic properties of superconductors of the second group,'' \textit{Soviet Physics JETP} 5, 1174--1182 (1957), for the type-II superconductivity vortex analogy.
\item Manton and Sutcliffe, \textit{Topological Solitons}, Cambridge University Press, 2004, for standard soliton, vortex, topological-charge, and moduli-space vocabulary.
\item Vilenkin and Shellard, \textit{Cosmic Strings and Other Topological Defects}, Cambridge University Press, 1994, for topological-defect language and cosmological context.
\item Bekenstein, Hawking, Page, and modern black-hole-information reviews, for any claim involving horizon entropy, thermality, reduced external states, or boundary information.
\end{itemize}

\section{What Is Claimed And What Remains Open}

This paper claims only the following:
\begin{itemize}
\item Field 01 can be organized into a coherent toy formalization chain;
\item closed phase nodes can be represented by winding;
\item a scalar modulus / VEV-profile layer can be modeled by a scalar profile and then given the Field 01 normal-retention interpretation;
\item gauge-like compensation reduces long-range phase mismatch;
\item a conserved current-like object follows from phase-shift symmetry;
\item memory can be formalized as an equivalence class;
\item boundary recording can be represented by a bulk-to-boundary map;
\item reduced external descriptions can be interpreted as limited access to boundary memory.
\item selected Field 01 terms can be tested against standard Abelian-Higgs, scalar-boundary, gauge-stiffness, overlap-bin, and Wilson-line diagnostics.
\end{itemize}

The following remain open:
\begin{itemize}
\item a precise fundamental action;
\item relation to the Standard Model;
\item derivation of spin, charge, and particle spectrum;
\item relation to the Higgs mechanism;
\item comparison with Abelian-Higgs and vortex literature;
\item construction of a boundary Hilbert space;
\item derivation of Hawking radiation or black-hole entropy;
\item testable physical predictions.
\item a controlled two-vortex interaction-energy comparison;
\item a canonical boundary/access rule for the connector class $\gamma_{\mathrm{class}}$.
\end{itemize}

\section{Conclusion}

The current mathematical development of Field 01 does not yet produce a completed physical theory. It does, however, produce a structured formalization path. Starting from phase circulation, one can define closed nodes, introduce a scalar modulus / VEV-profile layer, add the Field 01 normal-retention interpretation, add gauge-like compensation, obtain a current-like object, define memory as an equivalence class, map bulk memory to boundary memory, and interpret reduced external states as limited access to that boundary record. The diagnostic appendix adds a further discipline: key Field 01 phrases are tested against standard Abelian-Higgs benchmarks, scalar boundary-layer tension, gauge-compensation stiffness, overlap-motivated relation bins, and Wilson-line dressed phase bookkeeping.

The most important lesson is caution. Much of the mathematics resembles known vortex or Abelian-Higgs-like structures, and several first diagnostics reduce directly to standard objects. Field 01 must therefore be judged not by the mere presence of winding or gauge-like notation, but by whether its interpretation of scalar-profile retention, memory, relation classes, and boundary recording leads to a useful and eventually testable framework.

\end{document}